% GENERATED FILE. Edit canonical lessons, then run npm run generate:books.
% canonical-source-hash: fnv1a64:e2f52f743545795b
% canonical-lessons: LA-C424-compitum, LA-C424-bivium, LA-C424-flexus, LA-C424-anfractus, LA-C424-flecto

\chapter{Crossroads, Fork in the road, Bend, Winding, To bend}
\label{ch:la-crossroads-fork-in-the-road-bend-winding-to-bend}
\hlchaptermodality{424}

\begin{chapteropening}
\textbf{By the end of this chapter:} \emph{I can talk about where roads meet and turn: a crossroads, a fork, a bend, a winding stretch, and bending a way.}

Complete the last lesson of chapter 424: I can talk about where roads meet and turn: a crossroads, a fork, a bend, a winding stretch, and bending a way.
\end{chapteropening}

\section[compitum, compitī]{compitum, compitī — a crossroads, a place where roads meet}

\label{lesson:LA-C424-compitum}



\begin{quote}
Before the new one: say the Latin for a slope, then the Latin for to step.
\end{quote}

\subsection*{You'll want to know: compitum, compitī}
\textbf{compitum, compitī} — \textquotedblleft{}a crossroads, a place where roads meet\textquotedblright{}.

\begin{quote}
Ad compitum viātor stat et viam quaerit.
\end{quote}

\begin{quote}
\emph{At the crossroads the traveller stops and looks for the road.}
\end{quote}

Second declension, neuter: \emph{compitum, compitī n.} Small altars to the gods of the crossroads stood there.

\begin{cousinweb}
\emph{compitum} is usually explained as built from \emph{com-} 'together' and the root of \emph{petō} 'to make for', a place where ways come together. Lewis \& Short derive it from \emph{competō} 'to come together'; EtymDB lists it without a parent.
\end{cousinweb}

\begin{grammarlens}[title={What you've built}]
One more word for telling, explaining and describing at B1.
\end{grammarlens}

\subsection*{Your turn}
\begin{itemize}
\raggedright
  \item \emph{Say it:} \emph{compitum, compitī}
  \item \emph{Say it:} \emph{compitum, compitī}, once more
  \item \emph{Say it:} a sentence of your own with \emph{compitum, compitī}
  \item \emph{Recall:} say the Latin for a slope, then the Latin for to step, then say \emph{compitum, compitī} again
\end{itemize}

\begin{tcolorbox}[breakable,colback=teal!4,colframe=teal!35!black,title={Before you move on}]
What does compitum, compitī mean? (\textquotedblleft{}A crossroads, a place where roads meet\textquotedblright{}.) Say it once more.
\end{tcolorbox}
\section[bivium, biviī]{bivium, biviī — a fork in the road, a place where two roads meet}

\label{lesson:LA-C424-bivium}



\begin{quote}
Before the new one: say the Latin for to step, then the Latin for a crossroads.
\end{quote}

\subsection*{You'll want to know: bivium, biviī}
\textbf{bivium, biviī} — \textquotedblleft{}a fork in the road, a place where two roads meet\textquotedblright{}.

\begin{quote}
In biviō duae viae sunt: ūna ad mare, altera ad montēs.
\end{quote}

\begin{quote}
\emph{At the fork there are two roads: one to the sea, the other to the mountains.}
\end{quote}

Second declension, neuter: \emph{bivium, biviī n.} Compare \emph{trivium}, where three roads meet.

\begin{cousinweb}
\emph{bivium} is built on \emph{bi-} 'two' and \emph{via} 'road' (Lewis \& Short, under the adjective \emph{bivius} 'having two ways'). It is a native Latin compound.
\end{cousinweb}

\begin{grammarlens}[title={What you've built}]
One more word for telling, explaining and describing at B1.
\end{grammarlens}

\subsection*{Your turn}
\begin{itemize}
\raggedright
  \item \emph{Say it:} \emph{bivium, biviī}
  \item \emph{Say it:} \emph{bivium, biviī}, once more
  \item \emph{Say it:} a sentence of your own with \emph{bivium, biviī}
  \item \emph{Recall:} say the Latin for to step, then the Latin for a crossroads, then say \emph{bivium, biviī} again
\end{itemize}

\begin{tcolorbox}[breakable,colback=teal!4,colframe=teal!35!black,title={Before you move on}]
What does bivium, biviī mean? (\textquotedblleft{}A fork in the road, a place where two roads meet\textquotedblright{}.) Say it once more.
\end{tcolorbox}
\section[flexus, flexūs]{flexus, flexūs — a bend, a turn (in a road or river)}

\label{lesson:LA-C424-flexus}



\begin{quote}
Before the new one: say the Latin for a crossroads, then the Latin for a fork in the road.
\end{quote}

\subsection*{You'll want to know: flexus, flexūs}
\textbf{flexus, flexūs} — \textquotedblleft{}a bend, a turn (in a road or river)\textquotedblright{}.

\begin{quote}
Post flexum viae oppidum appāret.
\end{quote}

\begin{quote}
\emph{After the bend in the road the town comes into view.}
\end{quote}

Fourth declension, masculine: \emph{flexus, flexūs m.} English \emph{flex} and \emph{reflex} belong to the same family.

\begin{cousinweb}
\emph{flexus} is built on \emph{flectō} 'to bend' (EtymDB; Lewis \& Short), a fourth-declension noun of action made from the verb's stem \emph{flex-}. English \emph{flex} comes from the same verb.
\end{cousinweb}

\begin{grammarlens}[title={What you've built}]
One more word for telling, explaining and describing at B1.
\end{grammarlens}

\subsection*{Your turn}
\begin{itemize}
\raggedright
  \item \emph{Say it:} \emph{flexus, flexūs}
  \item \emph{Say it:} \emph{flexus, flexūs}, once more
  \item \emph{Say it:} a sentence of your own with \emph{flexus, flexūs}
  \item \emph{Recall:} say the Latin for a crossroads, then the Latin for a fork in the road, then say \emph{flexus, flexūs} again
\end{itemize}

\begin{tcolorbox}[breakable,colback=teal!4,colframe=teal!35!black,title={Before you move on}]
What does flexus, flexūs mean? (\textquotedblleft{}A bend, a turn (in a road or river)\textquotedblright{}.) Say it once more.
\end{tcolorbox}
\section[anfrāctus, anfrāctūs]{anfrāctus, anfrāctūs — a winding, a twisting stretch of road}

\label{lesson:LA-C424-anfractus}



\begin{quote}
Before the new one: say the Latin for a fork in the road, then the Latin for a bend.
\end{quote}

\subsection*{You'll want to know: anfrāctus, anfrāctūs}
\textbf{anfrāctus, anfrāctūs} — \textquotedblleft{}a winding, a twisting stretch of road\textquotedblright{}.

\begin{quote}
Via per montēs multīs anfrāctibus it.
\end{quote}

\begin{quote}
\emph{The road goes through the mountains with many windings.}
\end{quote}

Fourth declension, masculine: \emph{anfrāctus, anfrāctūs m.} Also a long, roundabout way of speaking.

\begin{cousinweb}
\emph{anfrāctus} is built on \emph{am(b)-} 'around' and \emph{frangō} 'to break' (Lewis \& Short; EtymDB lists \emph{frangō}, from PIE \emph{b\textsuperscript{h}reg-} 'to break'): a road that keeps breaking round in turns.
\end{cousinweb}

\begin{grammarlens}[title={What you've built}]
One more word for telling, explaining and describing at B1.
\end{grammarlens}

\subsection*{Your turn}
\begin{itemize}
\raggedright
  \item \emph{Say it:} \emph{anfrāctus, anfrāctūs}
  \item \emph{Say it:} \emph{anfrāctus, anfrāctūs}, once more
  \item \emph{Say it:} a sentence of your own with \emph{anfrāctus, anfrāctūs}
  \item \emph{Recall:} say the Latin for a fork in the road, then the Latin for a bend, then say \emph{anfrāctus, anfrāctūs} again
\end{itemize}

\begin{tcolorbox}[breakable,colback=teal!4,colframe=teal!35!black,title={Before you move on}]
What does anfrāctus, anfrāctūs mean? (\textquotedblleft{}A winding, a twisting stretch of road\textquotedblright{}.) Say it once more.
\end{tcolorbox}
\section[flectō, flectere]{flectō, flectere — to bend, to turn}

\label{lesson:LA-C424-flecto}



\begin{quote}
Before the new one: say the Latin for a bend, then the Latin for a winding.
\end{quote}

\subsection*{You'll want to know: flectō, flectere}
\textbf{flectō, flectere} — \textquotedblleft{}to bend, to turn\textquotedblright{}.

\begin{quote}
Agricola bovēs dextrā flectit, et plaustrum in viam angustam intrat.
\end{quote}

\begin{quote}
\emph{The farmer turns the oxen to the right, and the wagon goes into the narrow road.}
\end{quote}

\emph{flectō, flectere, flexī, flexum}. \emph{Iter flectere}, to change the direction of a journey. English \emph{flexible} and \emph{reflect} come from it.

\begin{cousinweb}
\emph{flectō} is inherited from Proto-Italic \emph{flektō} (EtymDB). English \emph{flexible}, \emph{reflect} and \emph{inflection} all come from it (Lewis \& Short).
\end{cousinweb}

\begin{grammarlens}[title={What you've built}]
One more word for telling, explaining and describing at B1.
\end{grammarlens}

\subsection*{Your turn}
\begin{itemize}
\raggedright
  \item \emph{Say it:} \emph{flectō, flectere}
  \item \emph{Say it:} \emph{flectō, flectere}, once more
  \item \emph{Say it:} a sentence of your own with \emph{flectō, flectere}
  \item \emph{Recall:} say the Latin for a bend, then the Latin for a winding, then say \emph{flectō, flectere} again
\end{itemize}

\begin{tcolorbox}[breakable,colback=teal!4,colframe=teal!35!black,title={Before you move on}]
What does flectō, flectere mean? (\textquotedblleft{}To bend, to turn\textquotedblright{}.) Say it once more.
\end{tcolorbox}
